% Original analytic teaching diagram, 2026-11-08.
\definecolor{VB}{HTML}{7188A3}\definecolor{VBF}{HTML}{E5EAF0}
\definecolor{VG}{HTML}{829C8B}\definecolor{VGF}{HTML}{E8EEE9}
\definecolor{VY}{HTML}{B4A16F}\definecolor{VYF}{HTML}{F1ECDD}
\definecolor{VR}{HTML}{A7787D}\definecolor{VRF}{HTML}{F0E5E6}
\definecolor{VP}{HTML}{9588A5}\definecolor{VPF}{HTML}{ECE8F0}
\definecolor{VI}{HTML}{3B4252}\definecolor{VM}{HTML}{6B7280}\definecolor{VL}{HTML}{4C4D4F}
\begin{tikzpicture}[x=1mm,y=1mm,>=Stealth,font=\figurefont\fontsize{8}{10}\selectfont,
lab/.style={text=VM,font=\figurefont\fontsize{7.5}{9}\selectfont}]
% (a) smooth curved domain with a local graph chart
\filldraw[fill=VBF,draw=VB,line width=.95pt] plot[smooth cycle,tension=.75] coordinates {(2,12) (11,3) (22,5) (29,16) (22,27) (9,25)};
\node[text=VI] at (13,15) {$\Omega$};
\draw[VY,line width=.95pt] (25.6,7.6) circle (6.6);
\fill[VI] (25.6,7.6) circle (.55);
\node[lab,anchor=north west] at (26.6,3.4) {$x_0$};
\draw[VL,line width=.65pt,->] (25.6,7.6) -- (20.9,12.2) node[pos=1.18,text=VI] {$x_d$};
\draw[VL,line width=.65pt,->] (25.6,7.6) -- (31.0,12.4) node[pos=1.2,text=VI] {$x'$};
\node[text=VI] at (15.5,-5) {(a)};
% (b) square: corner is still a Lipschitz graph
\begin{scope}[xshift=40mm]
\filldraw[fill=VBF,draw=VB,line width=.95pt] (3,3) rectangle (27,27);
\node[text=VI] at (13,16) {$\Omega$};
\draw[VY,line width=.95pt] (27,3) circle (6.6);
\fill[VI] (27,3) circle (.55);
\draw[VL,line width=.65pt,->] (27,3) -- (22.3,7.7) node[pos=1.2,text=VI] {$x_d$};
\node[text=VI] at (15,-5) {(b)};
\end{scope}
% (c) slit disk is not Lipschitz
\begin{scope}[xshift=80mm]
\filldraw[fill=VBF,draw=VB,line width=.95pt] (15,15) circle (13);
\draw[VR,line width=1.3pt] (15,15) -- (28,15);
\node[text=VI] at (9,20) {$\Omega$};
\draw[VY,line width=.95pt] (21,15) circle (4.6);
\fill[VI] (21,15) circle (.55);
\draw[VL,line width=.65pt,->] (21,16.1) -- (21,23.0);
\draw[VL,line width=.65pt,->] (21,13.9) -- (21,7.0);
\node[text=VI] at (15,-5) {(c)};
\end{scope}
\end{tikzpicture}
